\section{Universally catenary schemes}
\label{section-universally-catenary}

\noindent
Recall that a topological space $X$ is called {\it catenary} if
for every pair of irreducible closed subsets $T \subset T'$
there exist a maximal chain of irreducible closed subsets
$$
T = T_0 \subset T_1 \subset \ldots \subset T_e = T'
$$
and every such chain has the same length. See
Topology, Definition \ref{topology-definition-catenary}.
Recall that a scheme is catenary if its underlying topological space
is catenary. See Properties, Definition \ref{properties-definition-catenary}.

\begin{definition}
\label{definition-universally-catenary}
Let $S$ be a scheme. Assume $S$ is locally Noetherian.
We say $S$ is {\it universally catenary} if for every
morphism $X \to S$ locally of finite type the scheme $X$ is catenary.
\end{definition}

\noindent
This is a ``better'' notion than catenary as there exist Noetherian schemes
which are catenary but not universally catenary. See
Examples, Section \ref{examples-section-non-catenary-Noetherian-local}.
Many schemes are universally catenary, see
Lemma \ref{lemma-ubiquity-uc} below.

\medskip\noindent
Recall that a ring $A$ is called {\it catenary} if
for any pair of prime ideals $\mathfrak p \subset \mathfrak q$
there exists a maximal chain of primes
$$
\mathfrak p =
\mathfrak p_0 \subset \ldots \subset \mathfrak p_e
= \mathfrak q
$$
and all of these have the same length. See
Algebra, Definition \ref{algebra-definition-catenary}.
We have seen the relationship between catenary schemes and
catenary rings in Properties, Section \ref{properties-section-catenary}.
Recall that a ring $A$ is called {\it universally catenary} if
$A$ is Noetherian and for every finite type ring map $A \to B$
the ring $B$ is catenary. See
Algebra, Definition \ref{algebra-definition-universally-catenary}.
Many interesting rings which come up
in algebraic geometry satisfy this property.

\begin{lemma}
\label{lemma-universally-catenary-local}
Let $S$ be a locally Noetherian scheme. The following are equivalent
\begin{enumerate}
\item $S$ is universally catenary,
\item there exists an open covering of $S$ all of whose members are
universally catenary schemes,
\item for every affine open $\Spec(R) = U \subset S$ the ring
$R$ is universally catenary, and
\item there exists an affine open covering $S = \bigcup U_i$ such
that each $U_i$ is the spectrum of a universally catenary ring.
\end{enumerate}
Moreover, in this case any scheme locally of finite type over $S$
is universally catenary as well.
\end{lemma}

\begin{proof}
By Lemma \ref{lemma-immersion-locally-finite-type} an open immersion
is locally of finite type. A composition of morphisms locally of
finite type is locally of finite type
(Lemma \ref{lemma-composition-finite-type}). Thus it is clear that if $S$ is
universally catenary then any open and any scheme locally of finite
type over $S$ is universally catenary as well. This proves the final
statement of the lemma and that (1) implies (2).

\medskip\noindent
If $\Spec(R)$ is a universally catenary scheme, then every
scheme $\Spec(A)$ with $A$ a finite type $R$-algebra is
catenary. Hence all these rings $A$ are catenary by
Algebra, Lemma \ref{algebra-lemma-catenary}.
Thus $R$ is universally catenary. Combined with the remarks above we
conclude that (1) implies (3), and (2) implies (4). Of course
(3) implies (4) trivially.

\medskip\noindent
To finish the proof we show that (4) implies (1).
Assume (4) and let $X \to S$ be a morphism locally of finite type.
We can find an affine open covering $X = \bigcup V_j$ such that
each $V_j \to S$ maps into one of the $U_i$. By
Lemma \ref{lemma-locally-finite-type-characterize}
the induced ring map $\mathcal{O}(U_i) \to \mathcal{O}(V_j)$ is
of finite type. Hence $\mathcal{O}(V_j)$ is catenary. Hence
$X$ is catenary by Properties, Lemma \ref{properties-lemma-catenary-local}.
\end{proof}

\begin{lemma}
\label{lemma-universally-catenary-local-rings-universally-catenary}
Let $S$ be a locally Noetherian scheme.
The following are equivalent:
\begin{enumerate}
\item $S$ is universally catenary, and
\item all local rings $\mathcal{O}_{S, s}$ of $S$ are universally catenary.
\end{enumerate}
\end{lemma}

\begin{proof}
Assume that all local rings of $S$ are universally catenary.
Let $f : X \to S$ be locally of finite type.
We know that $X$ is catenary if and only if $\mathcal{O}_{X, x}$ is
catenary for all $x \in X$. If $f(x) = s$, then $\mathcal{O}_{X, x}$
is essentially of finite type over $\mathcal{O}_{S, s}$. Hence
$\mathcal{O}_{X, x}$ is catenary by the assumption that
$\mathcal{O}_{S, s}$ is universally catenary.

\medskip\noindent
Conversely, assume that $S$ is universally catenary. Let $s \in S$.
We may replace $S$ by an affine open neighbourhood of $s$ by
Lemma \ref{lemma-universally-catenary-local}. Say $S = \Spec(R)$
and $s$ corresponds to the prime ideal $\mathfrak p$. Any finite
type $R_{\mathfrak p}$-algebra $A'$ is of the form
$A_{\mathfrak p}$ for some finite type $R$-algebra $A$.
By assumption (and Lemma \ref{lemma-universally-catenary-local} if you like)
the ring $A$ is catenary, and hence $A'$ (a localization of $A$) is
catenary. Thus $R_{\mathfrak p}$ is universally catenary.
\end{proof}

\begin{lemma}
\label{lemma-catenary-check-irreducible}
Let $S$ be a locally Noetherian scheme. Then $S$ is universally catenary
if and only if the irreducible components of $S$ are universally catenary.
\end{lemma}

\begin{proof}
Omitted. For the affine case, please see
Algebra, Lemma \ref{algebra-lemma-catenary-check-irreducible}.
\end{proof}

\begin{lemma}
\label{lemma-ubiquity-uc}
The following types of schemes are universally catenary.
\begin{enumerate}
\item Any scheme locally of finite type over a field.
\item Any scheme locally of finite type over a Cohen-Macaulay scheme.
\item Any scheme locally of finite type over $\mathbf{Z}$.
\item Any scheme locally of finite type over a $1$-dimensional
Noetherian domain.
\item And so on.
\end{enumerate}
\end{lemma}

\begin{proof}
All of these follow from the fact that a
Cohen-Macaulay ring is universally catenary, see
Algebra, Lemma \ref{algebra-lemma-CM-ring-catenary}.
Also, use the last assertion of
Lemma \ref{lemma-universally-catenary-local}.
Some details omitted.
\end{proof}























